\section{行列による関係の表現}
\subsection{授業内容}
\subsubsection{科目の中でのこのコマの位置づけ}
線形代数についての基礎的なルールを習得する段階である。今回はより実践的・具体的に，データ行列をどのように線形代数で表現できるかを考える。
データ行列から分散共分散行列，相関行列へと形を変えることを学ぶ。
つづいて線形モデル，特に従属変数が明確な回帰分析モデルを行列で表現することを見，線形モデルとデザイン行列について考える。
さらに因子分析モデルを行列で表現することを考える。
行列で表現することで，一つの式の中に第一，第二定理の両方を含んだ形で表現できることを理解する。

\subsubsection{コマ主題細目}
\begin{description}
    \item[データの行列表現] データの行列表現を考え，平均ベクトル，平均偏差行列，分散共分散行列を算出するプロセスを学ぶ。標準偏差からなる対角行列をもちいて分散共分散行列を相関行列に変換する。
    \item[線形モデル] 回帰分析モデルを行列で表現すると重回帰分析も一気に表現できるところを見る。また逆行列の計算から回帰係数が算出できることを見る。
    \item[デザイン行列] 分散分析は回帰分析の特殊例であることを論じてきたが，行列で表現する時は分散分析のデザインを行列で表現することになる。自由度の概念にも目を向けながら，デザイン行列を確認する。
    \item[因子分析モデルの行列表現]行列のサイズに注意しながら，因子分析モデルを行列で表現するところを見る。
\end{description}
\subsubsection{キーワード}
\begin{itemize}
\item データの行列表現
\item 分散共分散行列，相関行列
\item デザイン行列
\end{itemize}
\subsection{授業情報}
\paragraph{コマの展開方法}
講義
\subsubsection{予習・復習課題}
\paragraph{予習}行列の掛け算がメインになってくるので，計算方法並びに計算結果のサイズを確認する方法を見ておこう。
\paragraph{復習}行列表現によって重回帰方程式が一つの形になることを確認する。平均値の差を見るために線形モデルが用いられることを確認する。また因子分析モデルを行列表現すると，一気に二つの定理が一つの式で表現できることを確認する。
